\titulo{Análise crítica de uma representação gráfica enganosa nas eleições de 2026}
\subtitulo{Pesquisa eleitoral AtlasIntel/Bloomberg e a circulação de uma imagem com dados adulterados}
\titulotraduzido{}

\autor{Fabricio Fernandes Santana}
\autorficha{SANTANA, Fabricio Fernandes}

\orientador{Não se aplica}
\coorientador{}

\curso{Administração Pública}
\programa{Administração Pública}
\area{}
\linhapesquisa{}

\cidade{Brasília}
\local{Brasília-DF}
\data{2026}
\datadefesa{}

\cutter{}
\numeropaginas{}
\cdd{}
\assuntos{1. Eleições 2026. 2. Gráficos eleitorais. 3. Desinformação. 4. Alfabetização de dados.}

\palavraschave{Eleições 2026. Pesquisa eleitoral. Visualização de dados. Desinformação. Alfabetização de dados.}
\keywords{2026 elections. Electoral polling. Data visualization. Disinformation. Data literacy.}
